%% kapa: exam=2 round=33 year=2022 period=3 subject=compensation_law subject_ko=감정평가 및 보상법규 seq=2 number=2 points=30
%% source: Q-Net 감정평가사 기출문제 (cst00309) artlSeq=5210838 file=raw/1008/5210838/2022년 제33회 감정평가사 2차시험 3교시 문제지(감정평가 및 보상법규).pdf
%% conversion: kordoc 4.24.51 -> tools/gen2.py (auto); review=proofread-gpt-astra; flags=none
\begin{longquestion}{【문제 2】}
국토교통부장관은 표준지로 선정된 A토지의 2022. 1. 1. 기준 공시지가를 1 m\textsuperscript{2}당 1,000만원으로 결정ㆍ공시하였다. 국토교통부장관은 A토지의 표준지공시지가를 산정함에 있어 「부동산 가격공시에 관한 법률」및 같은 법 시행령이 정하는 ‘토지의 일반적인 조사사항’ 이외에 국토교통부 훈령인 「표준지공시지가 조사ㆍ평가 기준」상 상업ㆍ업무용지 평가의 고려사항인 ‘배후지의 상태 및 고객의 질과 양’, ‘영업의 종류 및 경쟁의 상태’ 등을 추가적으로 고려하여 평가하였다. 甲은 X시에 상업용지인 B토지를 소유하고 있다. X시장은 A토지를 비교표준지로 선정하여 B토지에 대한 개별공시지가를 1 m\textsuperscript{2}당 1,541만원으로 결정ㆍ공시 후 이를 甲에게 통지하였다. 甲은 국토교통부장관이 A토지의 표준지공시지가를 단순히 행정청 내부에서만 효력을 가지는 국토교통부 훈령 형식의 「표준지공시지가 조사ㆍ평가 기준」이 정하는 바에 따라 평가함으로써 결과적으로 부동산 가격공시에 관한 법령이 직접 규정하지 않는 사항을 표준지공시지가 평가의 고려사항으로 삼은 것은 위법하다고 주장하고 있다. 다음 물음에 답하시오. (30점) \par
\medskip\noindent 물음 1) 「표준지공시지가 조사ㆍ평가 기준」의 법적 성질에 비추어 甲 주장의 타당성 여부를 설명하시오. (20점) \par
\medskip\noindent 물음 2) 甲은 「부동산 가격공시에 관한 법률」제11조에 따라 X시장에게 B토지의 개별공시지가에 대한 이의를 신청하였으나 기각되었다. 이 경우 甲이 기각결정에 불복하여 「행정심판법」상의 행정심판을 제기할 수 있는지 설명하시오. (10점) \par
\begin{framedtext}
참조 조문 \par
<부동산 가격공시에 관한 법률> \par
제11조(개별공시지가에 대한 이의신청) ① 개별공시지가에 이의가 있는 자는 그 결정ㆍ공시일부터 30일 이내에 서면으로 시장ㆍ군수 또는 구청장에게 이의를 신청할 수 있다. \par
<부동산 가격공시에 관한 법률 시행령> \par
제6조(표준지공시지가 조사ㆍ평가의 기준) ① 법 제3조제4항에 따라 국토교통부장관이 표준지공시지가를 조사ㆍ평가하는 경우 참작하여야 하는 사항의 기준은 다음 각 호와 같다. <각 호 생략> \par
② 표준지에 건물 또는 그 밖의 정착물이 있거나 지상권 또는 그 밖의 토지의 사용ㆍ수익을 제한하는 권리가 설정되어 있을 때에는 그 정착물 또는 권리가 존재하지 아니하는 것으로 보고 표준지공시지가를 평가하여야 한다. \par
③ 제1항 및 제2항에서 규정한 사항 외에 표준지공시지가의 조사ㆍ평가에 필요한 세부기준은 국토교통부장관이 정한다. \par
<표준지공시지가 조사ㆍ평가 기준> \par
제23조(상업ㆍ업무용지) 상업ㆍ업무용지(공공용지를 제외한다)는 토지의 일반적인 조사사항 이외에 다음 각 호의 사항 등을 고려하여 평가하되, 인근지역 또는 동일수급권 안의 유사지역에 있는 토지의 거래사례 등 가격자료를 활용하여 거래사례비교법으로 평가한다. <단서 생략> \par
1. 배후지의 상태 및 고객의 질과 양 \par
2. 영업의 종류 및 경쟁의 상태 \par
3.～6. <생략> \par
\end{framedtext}
\end{longquestion}
